% !TeX root = ../../main.tex

\section{Frameworks de Big Data: Apache Hadoop y Apache Flink}

\subsection{Introducción}

Los frameworks de Big Data permiten almacenar y procesar grandes cantidades
de información utilizando varias computadoras. Su objetivo es distribuir el
trabajo para obtener resultados en menos tiempo y continuar funcionando
aunque alguno de los equipos presente una falla.

Apache Hadoop y Apache Flink son proyectos de código abierto administrados por
la Apache Software Foundation. Sin embargo, tienen enfoques diferentes.
Hadoop se utiliza principalmente para almacenar información y ejecutar
procesos por lotes, mientras que Flink está especializado en el procesamiento
continuo de datos en tiempo real.

\section{Apache Hadoop}

\subsection{Descripción}

Apache Hadoop es un framework de código abierto diseñado para almacenar y
procesar grandes volúmenes de datos de manera distribuida. Divide la
información y el trabajo entre diferentes computadoras que forman un clúster.

Hadoop fue diseñado para trabajar con equipos de costo relativamente bajo y
para recuperarse automáticamente cuando uno de los nodos presenta una falla.
Puede utilizarse con datos estructurados, semiestructurados y no estructurados.

\subsection{Características}

Las principales características de Hadoop son:

\begin{itemize}
    \item Almacenamiento distribuido de grandes cantidades de información.
    \item Procesamiento paralelo en varias computadoras.
    \item Escalabilidad horizontal mediante la incorporación de nuevos nodos.
    \item Tolerancia a fallas por medio de replicación y recuperación de tareas.
    \item Uso de hardware comercial.
    \item Capacidad para procesar datos estructurados y no estructurados.
    \item Integración con numerosos proyectos de código abierto.
    \item Adecuado para procesos por lotes.
\end{itemize}

\subsection{Antecedentes e historia}

Hadoop tiene su origen en el proyecto de motor de búsqueda Nutch, desarrollado
por Doug Cutting y Mike Cafarella. Su diseño recibió influencia de los
documentos publicados por Google sobre el sistema de archivos Google File
System y el modelo de procesamiento MapReduce.

En 2006 se creó Hadoop como un proyecto independiente. Yahoo apoyó su
desarrollo y lo utilizó para procesar grandes cantidades de información de su
motor de búsqueda. Con el tiempo, Hadoop se convirtió en uno de los proyectos
más representativos del Big Data.

La versión 1.0 fue publicada en 2011 y anunciada oficialmente en enero de 2012.
Posteriormente, Hadoop 2 incorporó YARN para separar la administración de
recursos del procesamiento. Hadoop 3 añadió mejoras de rendimiento,
disponibilidad y aprovechamiento del almacenamiento. La versión estable más
reciente es Hadoop 3.5.0.

\subsection{Arquitectura}

Hadoop utiliza una arquitectura distribuida de tipo maestro-trabajador. Cada
computadora del clúster recibe el nombre de nodo.

La arquitectura puede entenderse mediante tres partes principales:

\begin{enumerate}
    \item \textbf{Almacenamiento con HDFS:} la información se divide en bloques
    y se distribuye entre diferentes DataNodes. El NameNode mantiene el registro
    de la ubicación de cada bloque.
    
    \item \textbf{Administración con YARN:} el ResourceManager distribuye los
    recursos del clúster, mientras que los NodeManagers controlan los recursos
    disponibles en cada computadora.
    
    \item \textbf{Procesamiento con MapReduce:} los trabajos se dividen en
    tareas pequeñas que pueden ejecutarse de forma paralela.
\end{enumerate}

En instalaciones de alta disponibilidad pueden existir NameNodes y
ResourceManagers de respaldo. Si el proceso principal falla, otro puede asumir
su función.

\subsection{Componentes}

Los cuatro componentes principales de Hadoop son:

\begin{itemize}
    \item \textbf{Hadoop Common:} contiene las bibliotecas y utilidades
    compartidas por los demás componentes.
    
    \item \textbf{HDFS:} sistema de archivos distribuido que almacena grandes
    cantidades de datos en varios nodos.
    
    \item \textbf{YARN:} administra la memoria, el procesador y los demás
    recursos disponibles en el clúster.
    
    \item \textbf{MapReduce:} modelo de programación para procesar grandes
    conjuntos de datos de manera paralela.
\end{itemize}

En HDFS también se encuentran los siguientes procesos:

\begin{itemize}
    \item \textbf{NameNode:} administra los nombres, directorios y metadatos.
    \item \textbf{DataNode:} guarda físicamente los bloques de información.
    \item \textbf{JournalNode:} ayuda a compartir los cambios de metadatos en
    configuraciones de alta disponibilidad.
\end{itemize}

En YARN participan:

\begin{itemize}
    \item \textbf{ResourceManager:} distribuye los recursos del clúster.
    \item \textbf{NodeManager:} controla los recursos de cada nodo.
    \item \textbf{ApplicationMaster:} coordina la ejecución de una aplicación.
    \item \textbf{Container:} representa los recursos asignados a una tarea.
\end{itemize}

\subsection{Herramientas y servicios}

Hadoop puede complementarse con diferentes herramientas:

\begin{itemize}
    \item \textbf{Apache Hive:} permite consultar datos mediante un lenguaje
    parecido a SQL.
    
    \item \textbf{Apache HBase:} base de datos NoSQL distribuida que utiliza
    HDFS como almacenamiento.
    
    \item \textbf{Apache Pig:} facilita la creación de procesos de
    transformación de datos.
    
    \item \textbf{Apache Spark:} permite ejecutar procesos analíticos sobre
    HDFS y utilizar los recursos de YARN.
    
    \item \textbf{Apache Oozie:} programa y coordina trabajos de Hadoop.
    
    \item \textbf{Apache ZooKeeper:} coordina aplicaciones distribuidas.
    
    \item \textbf{DistCp:} copia grandes cantidades de información entre
    clústeres y sistemas de archivos.
    
    \item \textbf{Hadoop Streaming:} permite escribir tareas MapReduce con
    lenguajes distintos de Java.
\end{itemize}

\subsection{Modelo de trabajo}

El modelo tradicional de trabajo de Hadoop es el procesamiento por lotes.
Primero se reúne una cantidad de datos, después se procesa y finalmente se
guarda el resultado.

Un trabajo MapReduce se desarrolla de la siguiente manera:

\begin{enumerate}
    \item Los datos de entrada se almacenan en HDFS.
    \item La información se divide en bloques.
    \item YARN asigna memoria y procesador a las tareas.
    \item Las tareas \textit{Map} leen y transforman los datos.
    \item Los resultados se agrupan y ordenan en la etapa
    \textit{shuffle and sort}.
    \item Las tareas \textit{Reduce} combinan los resultados.
    \item La información final se escribe nuevamente en HDFS.
\end{enumerate}

Este modelo es adecuado para procesar mucha información, pero normalmente no
produce respuestas inmediatas.

\subsection{Integración}

Hadoop puede integrarse con bases de datos, herramientas analíticas y
servicios en la nube. HDFS puede trabajar con Hive, HBase, Spark, Flink y otras
tecnologías del ecosistema Apache.

También existen conectores para sistemas de almacenamiento como Amazon S3,
Azure Blob Storage y Azure Data Lake Storage. Los resultados pueden enviarse a
bases de datos, aplicaciones empresariales o plataformas de visualización.

\subsection{Hardware y software}

Hadoop está escrito principalmente en Java y normalmente se instala en
servidores con sistemas operativos Linux o similares a Unix. Para realizar
pruebas puede ejecutarse en una sola computadora, pero en producción se
utiliza un clúster con varios nodos.

Los nodos trabajadores necesitan:

\begin{itemize}
    \item Procesadores con varios núcleos.
    \item Memoria RAM suficiente para ejecutar tareas.
    \item Discos de gran capacidad.
    \item Una conexión de red rápida y estable.
\end{itemize}

El NameNode necesita suficiente memoria porque conserva los metadatos del
sistema de archivos. Para instalaciones grandes se recomienda utilizar
servidores de respaldo, redes redundantes y varios discos por nodo.

\subsection{Seguridad}

Hadoop puede utilizar Kerberos para comprobar la identidad de usuarios y
servicios. HDFS también maneja propietarios, grupos, permisos y listas de
control de acceso.

Entre sus mecanismos de seguridad se encuentran:

\begin{itemize}
    \item Autenticación mediante Kerberos.
    \item Permisos para archivos y directorios de HDFS.
    \item Tokens de delegación para aplicaciones.
    \item Cifrado de información almacenada.
    \item Cifrado de la comunicación entre nodos.
    \item Uso de HTTPS en las interfaces web.
    \item Hadoop KMS para administrar claves de cifrado.
    \item Separación de cuentas para HDFS, YARN y MapReduce.
\end{itemize}

Estas funciones deben configurarse correctamente. Una instalación básica de
Hadoop no debe exponerse directamente a Internet.

\subsection{Costos}

Apache Hadoop se distribuye bajo la licencia Apache 2.0, por lo que no requiere
el pago de una licencia para utilizar su código. Sin embargo, implementar un
clúster sí genera costos.

Los principales gastos son:

\begin{itemize}
    \item Compra o renta de servidores.
    \item Discos para guardar los datos y sus copias.
    \item Consumo de electricidad y refrigeración.
    \item Redes y equipos de respaldo.
    \item Personal para instalar, asegurar y administrar el clúster.
    \item Soporte o distribuciones comerciales.
    \item Servicios de nube cuando el clúster es rentado.
\end{itemize}

El costo final depende de la cantidad de datos, el número de nodos y el tiempo
que permanezcan activos.

\subsection{Ventajas y desventajas}

\textbf{Ventajas:}

\begin{itemize}
    \item Es software de código abierto.
    \item Puede almacenar cantidades muy grandes de información.
    \item Permite aumentar la capacidad agregando más nodos.
    \item Tolera fallas de discos y computadoras.
    \item Tiene un ecosistema amplio de herramientas.
    \item Es adecuado para procesamiento histórico por lotes.
\end{itemize}

\textbf{Desventajas:}

\begin{itemize}
    \item Su instalación y administración pueden ser complicadas.
    \item MapReduce puede ser lento para consultas interactivas.
    \item La replicación de datos consume espacio adicional.
    \item Requiere personal con conocimientos especializados.
    \item No es la mejor opción para respuestas en tiempo real.
    \item Los archivos pequeños pueden generar problemas en el NameNode.
\end{itemize}

\subsection{Casos de uso}

Hadoop puede utilizarse en los siguientes casos:

\begin{itemize}
    \item Almacenamiento de registros de aplicaciones.
    \item Construcción de lagos de datos.
    \item Análisis histórico de ventas.
    \item Procesamiento de documentos, imágenes y archivos.
    \item Creación de índices para motores de búsqueda.
    \item Preparación de datos para modelos de aprendizaje automático.
    \item Análisis de información producida por sensores.
    \item Copias y conservación de grandes conjuntos de datos.
\end{itemize}

\subsection{Caso de éxito}

Yahoo fue una de las primeras organizaciones que utilizó Hadoop a gran escala.
En 2012, la Apache Software Foundation informó que su infraestructura Hadoop
contaba con más de 42,000 nodos. La plataforma procesaba información
relacionada con búsquedas, publicidad y personalización del contenido para
cientos de millones de usuarios.

Este caso ayudó a demostrar que Hadoop podía utilizarse para almacenar y
procesar grandes volúmenes de información en ambientes reales.

\subsection{Video de ejemplo}

El siguiente video explica de forma sencilla qué es Apache Hadoop, sus
componentes y algunos casos de uso:

\url{https://www.youtube.com/watch?v=JWX5Inb--ig}


\section{Apache Flink}

\subsection{Descripción}

Apache Flink es un framework y motor distribuido para procesar flujos de datos.
Puede trabajar con flujos ilimitados que reciben información continuamente y
con conjuntos de datos limitados que ya se encuentran almacenados.

Su principal especialidad es el procesamiento en tiempo real. Flink puede
analizar cada evento mientras se produce, conservar el estado de una
aplicación y generar resultados con baja latencia.

\subsection{Características}

Las principales características de Flink son:

\begin{itemize}
    \item Procesamiento de datos en tiempo real.
    \item Compatibilidad con flujos limitados e ilimitados.
    \item Procesamiento con estado.
    \item Manejo de tiempo de evento y marcas de agua.
    \item Procesamiento paralelo y distribuido.
    \item Baja latencia y alto rendimiento.
    \item Recuperación mediante checkpoints.
    \item Garantías de procesamiento exactamente una vez.
    \item APIs para Java, Python, SQL y Table API.
    \item Integración con diferentes sistemas de almacenamiento y mensajería.
\end{itemize}

\subsection{Antecedentes e historia}

Flink tiene su origen en el proyecto de investigación Stratosphere, iniciado
en 2009 en la Universidad Técnica de Berlín y otras instituciones europeas.

En 2014, el proyecto ingresó a la Apache Software Foundation y cambió su nombre
a Apache Flink. En diciembre del mismo año fue aprobado como proyecto de nivel
superior de Apache, y el anuncio oficial se realizó en enero de 2015.

Con el paso del tiempo, Flink evolucionó para convertirse en una plataforma
especializada en procesamiento continuo, SQL, análisis por lotes y
aplicaciones dirigidas por eventos. La versión estable más reciente es Apache
Flink 2.3.0.

\subsection{Arquitectura}

Flink utiliza una arquitectura maestro-trabajador. Sus dos procesos
principales son el JobManager y los TaskManagers.

\begin{itemize}
    \item \textbf{JobManager:} coordina la ejecución de los programas,
    distribuye tareas y administra la recuperación.
    
    \item \textbf{TaskManager:} ejecuta las tareas y procesa los datos.
    
    \item \textbf{Task Slots:} representan las unidades de recursos que un
    TaskManager puede asignar a las tareas.
\end{itemize}

El JobManager contiene otros componentes:

\begin{itemize}
    \item \textbf{ResourceManager:} administra los recursos y task slots.
    \item \textbf{Dispatcher:} recibe las aplicaciones y proporciona la
    interfaz web.
    \item \textbf{JobMaster:} coordina la ejecución de un trabajo específico.
\end{itemize}

Una aplicación puede ejecutarse en un clúster compartido de sesión o en un
clúster dedicado. Flink también puede trabajar de forma independiente o sobre
YARN y Kubernetes.

\subsection{Componentes}

Los componentes principales de Apache Flink son:

\begin{itemize}
    \item \textbf{Flink Runtime:} motor que ejecuta las tareas distribuidas.
    
    \item \textbf{DataStream API:} permite crear aplicaciones de procesamiento
    continuo.
    
    \item \textbf{Table API:} ofrece operaciones sobre datos organizados como
    tablas.
    
    \item \textbf{Flink SQL:} permite consultar flujos y datos almacenados con
    instrucciones SQL.
    
    \item \textbf{State Backends:} almacenan el estado de las aplicaciones.
    
    \item \textbf{Checkpoints:} guardan copias automáticas del estado para
    recuperar una aplicación.
    
    \item \textbf{Savepoints:} permiten detener, actualizar o mover una
    aplicación conservando su estado.
    
    \item \textbf{Connectors:} comunican Flink con fuentes y destinos externos.
    
    \item \textbf{CEP:} biblioteca para detectar patrones complejos de eventos.
\end{itemize}

\subsection{Herramientas y servicios}

Flink ofrece diferentes herramientas para desarrollar y operar aplicaciones:

\begin{itemize}
    \item \textbf{Flink SQL Client:} ejecuta consultas SQL sobre datos.
    \item \textbf{Web Dashboard:} muestra trabajos, tareas, recursos y errores.
    \item \textbf{Command Line Interface:} permite enviar y administrar trabajos.
    \item \textbf{REST API:} facilita la administración mediante aplicaciones.
    \item \textbf{Flink CDC:} captura cambios realizados en bases de datos.
    \item \textbf{Flink Kubernetes Operator:} administra aplicaciones de Flink
    en Kubernetes.
    \item \textbf{PyFlink:} permite desarrollar aplicaciones utilizando Python.
\end{itemize}

También existen servicios administrados, como Amazon Managed Service for
Apache Flink y ofertas comerciales basadas en Flink.

\subsection{Modelo de trabajo}

Flink representa una aplicación mediante un flujo de datos formado por
fuentes, transformaciones y destinos.

El proceso general es el siguiente:

\begin{enumerate}
    \item Una fuente recibe información desde Kafka, archivos, sensores o bases
    de datos.
    
    \item Los eventos se distribuyen entre operadores paralelos.
    
    \item Los operadores filtran, agrupan, relacionan o transforman la
    información.
    
    \item La aplicación conserva el estado necesario para procesar eventos
    posteriores.
    
    \item Los checkpoints guardan copias del estado.
    
    \item Los resultados se envían a una base de datos, otro flujo, un archivo
    o un tablero.
\end{enumerate}

Flink puede procesar eventos usando el momento en que fueron creados, conocido
como tiempo de evento. Las marcas de agua ayudan a manejar eventos que llegan
con retraso.

\subsection{Integración}

Flink cuenta con conectores para sistemas como:

\begin{itemize}
    \item Apache Kafka y Apache Pulsar.
    \item Amazon Kinesis y Google Cloud Pub/Sub.
    \item RabbitMQ.
    \item HDFS y sistemas de archivos.
    \item Apache Hive y Apache HBase.
    \item Cassandra y MongoDB.
    \item Elasticsearch y OpenSearch.
    \item Bases de datos compatibles con JDBC.
    \item Amazon S3 y otros almacenamientos de objetos.
\end{itemize}

También puede ejecutarse sobre Hadoop YARN, Kubernetes, Docker o un clúster
independiente. Los formatos compatibles incluyen CSV, JSON, Avro, Parquet,
ORC y Protobuf.

\subsection{Hardware y software}

Flink puede instalarse en una computadora para realizar pruebas o en un
clúster para trabajar en producción. La versión 2.3 admite Java 11, 17 o 21.
PyFlink también requiere Python 3.9 o una versión compatible posterior.

El hardware depende de la cantidad y velocidad de los eventos. Generalmente
se necesita:

\begin{itemize}
    \item Procesadores con varios núcleos.
    \item Memoria RAM para mantener el estado.
    \item Almacenamiento para checkpoints y savepoints.
    \item Una red rápida entre TaskManagers.
    \item Nodos de respaldo para alta disponibilidad.
\end{itemize}

Para aplicaciones pequeñas puede utilizarse Docker. Las aplicaciones de gran
escala suelen desplegarse en Kubernetes, YARN o servicios cloud.

\subsection{Seguridad}

Flink puede proteger sus comunicaciones internas y externas mediante TLS o
SSL. También puede autenticarse con sistemas externos utilizando Kerberos y
tokens de delegación.

Sus medidas de seguridad incluyen:

\begin{itemize}
    \item Cifrado de la comunicación entre JobManager y TaskManagers.
    \item Protección de la interfaz REST mediante HTTPS.
    \item Autenticación mutua con certificados.
    \item Integración con Kerberos.
    \item Tokens para acceder a HDFS, Kafka y otros servicios.
    \item Uso de proxies como NGINX o Envoy para proteger la interfaz REST.
    \item Permisos y controles proporcionados por Kubernetes o YARN.
\end{itemize}

Varias opciones de seguridad no se encuentran activadas en una instalación
básica, por lo que deben configurarse antes de utilizar Flink en producción.

\subsection{Costos}

Apache Flink también utiliza la licencia Apache 2.0 y puede descargarse sin
pagar una licencia. Los costos aparecen al proporcionar la infraestructura
necesaria.

Los gastos principales son:

\begin{itemize}
    \item Servidores o máquinas virtuales.
    \item Memoria para mantener el estado.
    \item Almacenamiento para checkpoints.
    \item Transferencia de datos por la red.
    \item Administración y supervisión de las aplicaciones.
    \item Servicios administrados contratados con un proveedor cloud.
\end{itemize}

Un clúster compartido puede aprovechar mejor los recursos, mientras que un
clúster dedicado ofrece más aislamiento. La configuración debe ajustarse para
evitar mantener recursos sin utilizar.

\subsection{Ventajas y desventajas}

\textbf{Ventajas:}

\begin{itemize}
    \item Procesa información con baja latencia.
    \item Trabaja con datos continuos y datos por lotes.
    \item Proporciona manejo avanzado del estado.
    \item Cuenta con recuperación automática mediante checkpoints.
    \item Ofrece APIs de diferentes niveles.
    \item Se integra con Kafka, Hadoop, bases de datos y servicios cloud.
    \item Puede escalar a una gran cantidad de procesadores y datos.
\end{itemize}

\textbf{Desventajas:}

\begin{itemize}
    \item Su configuración puede ser complicada.
    \item Las aplicaciones con mucho estado requieren bastante memoria.
    \item Es necesario comprender conceptos como tiempo de evento y marcas de
    agua.
    \item Los checkpoints pueden consumir almacenamiento.
    \item La supervisión de aplicaciones continuas requiere personal
    especializado.
    \item Los conectores deben ser compatibles con la versión instalada.
\end{itemize}

\subsection{Casos de uso}

Apache Flink puede emplearse en:

\begin{itemize}
    \item Detección de fraudes en tiempo real.
    \item Seguimiento de transacciones financieras.
    \item Procesamiento de sensores y dispositivos IoT.
    \item Recomendaciones de productos o contenido.
    \item Supervisión de redes y aplicaciones.
    \item Detección de comportamientos anormales.
    \item Procesos ETL continuos.
    \item Actualización de índices de búsqueda.
    \item Análisis de registros y eventos.
    \item Procesamiento por lotes de información histórica.
\end{itemize}

\subsection{Caso de éxito}

Alibaba utiliza una versión de Flink llamada Blink para procesar información
en tiempo real. La empresa analiza búsquedas, clics y actividad de los usuarios
para mejorar la clasificación de productos y generar recomendaciones.

Este caso demuestra que Flink puede procesar flujos de gran tamaño y producir
resultados rápidamente en una plataforma de comercio electrónico.

\subsection{Video de ejemplo}

El siguiente video explica el funcionamiento de Apache Flink y su integración
con plataformas de eventos como Apache Kafka:

\url{https://www.youtube.com/watch?v=X1bhZZgAlCI}


\section{Comparación entre Hadoop y Flink}

\begin{itemize}
    \item \textbf{Tipo de procesamiento:} Hadoop se enfoca principalmente en
    lotes, mientras que Flink se especializa en flujos continuos.
    
    \item \textbf{Velocidad:} Hadoop puede tardar minutos u horas en completar
    un trabajo. Flink puede entregar resultados en segundos o milisegundos.
    
    \item \textbf{Almacenamiento:} Hadoop incluye HDFS. Flink necesita utilizar
    un sistema externo para guardar permanentemente los datos.
    
    \item \textbf{Estado:} Flink cuenta con herramientas avanzadas para
    conservar el estado de aplicaciones continuas.
    
    \item \textbf{Recuperación:} Hadoop repite las tareas que fallan. Flink
    recupera el estado mediante checkpoints.
    
    \item \textbf{Uso recomendado:} Hadoop resulta adecuado para información
    histórica y lagos de datos. Flink es conveniente para alertas, monitoreo y
    análisis en tiempo real.
\end{itemize}

Los dos frameworks pueden complementarse. Por ejemplo, Hadoop puede almacenar
la información histórica en HDFS y Flink puede leerla junto con eventos en
tiempo real.

\section{Conclusión}

Apache Hadoop y Apache Flink son herramientas importantes en el procesamiento
de Big Data. Hadoop proporciona almacenamiento distribuido y procesamiento
por lotes confiable. Flink ofrece procesamiento continuo, baja latencia y
manejo avanzado del estado.

La elección depende del problema que se desea resolver. Hadoop es apropiado
cuando se necesita conservar y procesar grandes cantidades de información
histórica. Flink es una mejor alternativa cuando los datos deben analizarse
mientras se producen. En algunas arquitecturas se pueden utilizar ambos para
combinar análisis histórico y procesamiento en tiempo real.